\documentclass[UTF8,a4paper,12pt]{ctexart}
\usepackage{geometry,amsmath,booktabs,hyperref}
\geometry{left=2.5cm,right=2.5cm,top=2.5cm,bottom=2.5cm}
\title{高性能芯片热管理系统问题三：多目标优化}
\author{数学建模项目组}\date{\today}
\begin{document}\maketitle
\section{问题描述}
在问题二得到的三个二次响应面代理模型基础上，寻求热阻、压降和温度非均匀性均较小的结构参数组合。令 $x_1,x_2,x_3$ 分别为针肋宽度比、歧管深高比和针肋排数，三个预测指标记为 $\hat R(x)$、$\hat P(x)$ 和 $\hat U(x)$。
\section{多目标优化模型}
三个目标均为极小化问题：
\begin{equation}
\min_{x\in\Omega}\ F(x)=\big(\hat R(x),\hat P(x),\hat U(x)\big),
\end{equation}
其中约束域取附件2实际试验水平的笛卡尔积
\[
\Omega=\{0,0.10,0.15,0.20,0.30\}\times\{3.0,3.5,4.0,4.5\}\times\{2,4,6,8,10\}.
\]
这样可以避免代理模型在样本范围外外推，并保证设计参数可直接实施。

由于三个指标量纲和数值范围不同，采用附件2样本极值进行极差归一化：
\begin{equation}
 z_j(x)=\frac{\hat y_j(x)-y_j^{\min}}{y_j^{\max}-y_j^{\min}},\qquad j=1,2,3,
\end{equation}
其中 $y_1=R,y_2=P,y_3=U$。采用等权综合评分作为综合最优准则：
\begin{equation}
 S(x)=\frac{z_1(x)+z_2(x)+z_3(x)}{3},\qquad \min_{x\in\Omega}S(x).
\end{equation}
等权意味着在题目未指定应用偏好的情况下，三项性能同等重要。为避免综合评分掩盖单项极差，另计算 Pareto 非支配集：若不存在 $x'$ 使三个指标均不大于 $x$ 且至少一项严格更小，则 $x$ 为非支配解。

\section{优化结果}
\begin{table}[h]\centering\caption{综合最优设计方案}
\begin{tabular}{cccccc}\toprule
$x_1$ & $x_2$ & $x_3$ & $\hat R$ & $\hat P$ & $\hat U$\\\midrule
0.20&4.0&2&0.755886&0.083295&0.784046\\
\bottomrule\end{tabular}\end{table}

该点归一化指标为
\[
(z_R,z_P,z_U)=(0.65491,0.05188,0.10128),
\]
综合评分
\[
S^*=0.269355.
\]
因此，在等权且参数限定为附件2试验水平的条件下，推荐
\[
\boxed{x_1=0.20,\qquad x_2=4.0,\qquad x_3=2.}
\]

需要指出，样本64 $(0.20,4.5,4)$ 的实测指标也较均衡，而代理模型在 $(0.20,4.0,2)$ 处给出的压降和温度非均匀性更优。两者差异来自代理模型的平滑拟合，最终实施前应对推荐点进行一次高保真仿真或实验复核。

\section{结果解释}
最优方案采用中等针肋宽度比，避免 $x_1=0.30$ 带来的明显阻塞；歧管深高比取4.0，在降低压降和维持换热之间折中；针肋排数取2，牺牲少量热阻以显著降低压降并获得较低温度非均匀性。该结果体现了三目标之间的冲突，不能简单将热阻单项最小作为设计依据。

\section{结论}
本文建立了基于二次响应面的离散多目标优化模型，并以极差归一化、等权综合评分和 Pareto 非支配性作为合理准则。得到的综合最优候选为 $(x_1,x_2,x_3)=(0.20,4.0,2)$，预测 $(R,P,U)=(0.755886,0.083295,0.784046)$。该方案应作为问题四权重敏感性和鲁棒优化的基准点。
\end{document}
在25?个离散组合（实际为 $5\times4\times5=100$ 个组合）上逐点计算问题二代理模型，得到等权评分最小且属于 Pareto 集的方案：
